% ============================================================================
% RAPPORT DE STAGE 5A - PROJET DE FIN D'ÉTUDES
% BEN SAID Nassim - INSA Hauts-de-France - Cycle Ingénieur ICY
% Stage : Bahgat IT Expert (Dubaï) - Plateforme de Business Intelligence
%         pilotée par l'Intelligence Artificielle
%
% >>> OVERLEAF : compiler avec pdfLaTeX. Créer un dossier "images/" et y
% >>> déposer vos captures. Chaque emplacement d'image est signalé par un
% >>> cadre gris "IMAGE À INSÉRER" + une ligne \includegraphics commentée
% >>> juste au-dessus : décommentez-la et supprimez le placeholder.
%
% >>> Les chiffres marqués "% MAJ" proviennent de votre exécution du
% >>> pipeline : vérifiez-les après votre DERNIÈRE exécution (page
% >>> "Data Pipeline" de l'application ou sortie console).
% ============================================================================
\documentclass[12pt,a4paper]{report}

\usepackage[utf8]{inputenc}
\usepackage[T1]{fontenc}
\usepackage[french]{babel}
\usepackage[a4paper,margin=2.5cm,headheight=15pt]{geometry}
\usepackage{graphicx}
\graphicspath{{assets_rs/}}
\usepackage{xcolor}
\usepackage{booktabs}
\usepackage{array}
\usepackage{amsmath,amssymb}
\usepackage{caption}
\usepackage{setspace}
\usepackage{enumitem}
\usepackage{fancyhdr}
\usepackage{titlesec}
\usepackage{listings}
\usepackage{microtype}
\usepackage[hidelinks]{hyperref}

% ---------------------------------------------------------------- couleurs
\definecolor{bipblue}{HTML}{1C5CAB}
\definecolor{bipgray}{HTML}{52514E}
\definecolor{phbg}{HTML}{F0EFEC}

% ---------------------------------------------------------------- en-têtes
\pagestyle{fancy}
\fancyhf{}
\fancyhead[L]{\small\nouppercase{\leftmark}}
\fancyhead[R]{\small BEN SAID Nassim}
\fancyfoot[C]{\thepage}
\renewcommand{\headrulewidth}{0.4pt}

% ---------------------------------------------------------------- titres
\titleformat{\chapter}[block]
  {\normalfont\huge\bfseries\color{bipblue}}{\thechapter.}{0.6em}{}
\titlespacing*{\chapter}{0pt}{-20pt}{25pt}

% ---------------------------------------------------------------- interligne
\onehalfspacing
\setlength{\parskip}{4pt}

% ---------------------------------------------------------------- listings
\lstset{
  basicstyle=\ttfamily\footnotesize,
  keywordstyle=\color{bipblue}\bfseries,
  commentstyle=\color{bipgray}\itshape,
  stringstyle=\color{teal},
  frame=single, rulecolor=\color{gray!40},
  breaklines=true, showstringspaces=false,
  captionpos=b, tabsize=4,
  literate={é}{{\'e}}1 {è}{{\`e}}1 {à}{{\`a}}1 {ç}{{\c c}}1 {ù}{{\`u}}1
           {ê}{{\^e}}1 {î}{{\^\i}}1 {ô}{{\^o}}1 {û}{{\^u}}1 {É}{{\'E}}1
}

% ------------------------------------------------- placeholder d'image
% Usage : \imageplaceholder{hauteur}{description de la capture à insérer}
\newcommand{\imageplaceholder}[2]{%
  \begin{center}
    \fcolorbox{gray!60}{phbg}{%
      \parbox[c][#1][c]{0.88\textwidth}{%
        \centering
        {\color{bipgray}\bfseries\large IMAGE À INSÉRER}\\[6pt]
        {\color{bipgray}\small #2}%
      }}%
  \end{center}}

% ============================================================================
\begin{document}

% ============================================================================
% PAGE DE GARDE
% ============================================================================
\begin{titlepage}
  \centering

  % --- logos : INSA HdF à gauche, Bahgat IT Expert à droite -----------------
   \includegraphics[height=2cm]{assets_rs/Insa-hautsdefrance-logo.png}\hfill
   \includegraphics[height=2cm]{assets_rs/AB_consulting.png}
  

  \vspace{2.2cm}

  {\large BEN SAID Nassim -- 5A ICY\par} 
  \vspace{0.4cm}
  {\large Rapport de Stage de Fin d'Études\\
   5\ieme{} année du Cycle Ingénieur ICY -- INSA Hauts-de-France\par}
  \vspace{0.3cm}
  {\large 1\ier{} avril 2026 -- 30 septembre 2026\par}

  \vspace{1.6cm}

  {\LARGE\bfseries\color{bipblue}
   CONCEPTION ET DÉVELOPPEMENT D'UNE\\
   PLATEFORME DE BUSINESS INTELLIGENCE\\
   
   PILOTÉE PAR L'INTELLIGENCE ARTIFICIELLE\par}

  \vspace{1.6cm}

  {\large
   Tuteur stage : \textit{BAHGAT Ahmed} \,|\, Mail : \textit{abahgat@bahgatexpert.com}\\ 
   Tuteur école : \textit{NIAR Smail} \,|\, Mail : \textit{Smail.Niar@uphf.fr}\par} 

  \vspace{1.4cm}

  {\large
   Société : \textbf{Bahgat IT Expert} -- Consulting \& Arbitration\\
   Adresse : Centurion Star Building, Port Said Avenue, Deira, Dubaï, EAU\\
   Secteur : Conseil en technologies de l'information\par}

  \vfill

  % --- éventuelle illustration de couverture --------------------------------
  % \includegraphics[width=0.6\textwidth]{images/couverture_dashboard.png}

\end{titlepage}

% ============================================================================
% REMERCIEMENTS
% ============================================================================
\chapter*{Remerciements}
\addcontentsline{toc}{chapter}{Remerciements}

Je tiens tout d'abord à remercier l'ensemble des équipes de Bahgat IT Expert
pour leur accueil et leur disponibilité tout au long de ces six mois de stage
effectués à distance.

Mes remerciements s'adressent en particulier à \textit{BAHGAT Ahmed}, mon tuteur de stage, pour la confiance qu'il m'a accordée en me
confiant un projet aussi structurant pour la practice Data \& Business
Intelligence du cabinet, ainsi que pour ses conseils avisés, tant techniques que
méthodologiques. 

Je remercie également \textit{NIAR Smail}, mon tuteur académique à
l'INSA Hauts-de-France, pour son suivi régulier et ses recommandations dans la
préparation de ce rapport et de la soutenance.

Enfin, je souhaite exprimer ma gratitude envers l'ensemble du corps enseignant
du cycle ingénieur ICY, dont les enseignements en science des données, en
développement logiciel et en gestion de projet ont constitué le socle sur
lequel ce projet de fin d'études a pu être construit.

% ============================================================================
% RÉSUMÉ / ABSTRACT
% ============================================================================
\chapter*{Résumé}
\addcontentsline{toc}{chapter}{Résumé}

Ce rapport présente le projet de fin d'études que j'ai réalisé au sein du
cabinet de conseil Bahgat IT Expert (Dubaï, Émirats Arabes Unis) : la
conception et le développement d'une plateforme de \textit{Business
Intelligence} pilotée par l'intelligence artificielle, destinée à transformer
des données métier brutes et multi-sources en indicateurs actionnables pour les
décideurs.

La plateforme couvre l'intégralité de la chaîne de valeur de la donnée : un
pipeline automatisé d'ingestion et de fiabilisation des données (ETL), un
entrepôt structuré en schéma en étoile, trois modules de machine learning
(prévision de chiffre d'affaires à 90 jours, détection d'anomalies
opérationnelles, segmentation client RFM avec score d'attrition), et deux
couches de restitution complémentaires : une application web interactive
développée en Flask et un rapport Power BI multi-pages.

Les données des clients du cabinet étant strictement confidentielles, la
démonstration repose sur un jeu de données synthétique haute-fidélité
(environ un million de lignes de ventes sur trois ans) reproduisant les
propriétés statistiques du retail du Golfe -- saisonnalités du Ramadan et des
grands événements commerciaux, croissance du canal e-commerce, promotions --
et dans lequel des incidents et des défauts de qualité connus ont été
volontairement injectés afin d'évaluer objectivement chaque brique de la
plateforme.

\vspace{0.5em}
\noindent\textbf{Mots-clés} : Business Intelligence, Machine Learning,
prévision, détection d'anomalies, segmentation RFM, ETL, entrepôt de données,
Flask, Power BI, Python.

\section*{Abstract}
This report presents my final-year engineering project at Bahgat IT Expert
(Dubai, UAE): the design and development of an AI-driven Business Intelligence
platform that turns raw, multi-source business data into actionable insights
for decision-makers. The platform covers the full data value chain: automated
data pipelines with quality remediation, a star-schema data warehouse, three
machine-learning modules (90-day revenue forecasting, anomaly detection, RFM
customer segmentation with churn scoring) and two reporting layers -- an
interactive Flask web application and a multi-page Power BI report. As client
data is confidential, the demonstration relies on a high-fidelity synthetic
dataset ($\sim$1M sales lines over three years) with known injected incidents,
enabling an objective evaluation of every component.

\vspace{0.5em}
\noindent\textbf{Keywords}: Business Intelligence, Machine Learning,
forecasting, anomaly detection, RFM segmentation, ETL, data warehouse, Flask,
Power BI, Python.

% ============================================================================
% TABLES
% ============================================================================
\tableofcontents
\listoffigures
\addcontentsline{toc}{chapter}{Table des figures}

% ============================================================================
% GLOSSAIRE
% ============================================================================
\chapter*{Glossaire et acronymes}
\addcontentsline{toc}{chapter}{Glossaire et acronymes}

\begin{description}[style=nextline,leftmargin=3.2cm,labelwidth=2.8cm]
  \item[BI] \textit{Business Intelligence} : ensemble des méthodes et outils
    permettant de transformer les données d'une organisation en information
    utile à la décision.
  \item[ETL] \textit{Extract, Transform, Load} : processus d'extraction, de
    transformation et de chargement des données vers un entrepôt.
  \item[KPI] \textit{Key Performance Indicator} : indicateur clé de
    performance.
  \item[Schéma en étoile] Modèle d'entrepôt de données (Kimball) séparant les
    tables de faits (mesures) des dimensions (axes d'analyse).
  \item[MAPE] \textit{Mean Absolute Percentage Error} : erreur absolue moyenne
    en pourcentage, métrique d'évaluation des prévisions.
  \item[RMSE] \textit{Root Mean Squared Error} : racine de l'erreur
    quadratique moyenne.
  \item[Holt-Winters] Méthode de lissage exponentiel triple modélisant niveau,
    tendance et saisonnalité d'une série temporelle.
  \item[Gradient Boosting] Méthode d'apprentissage supervisé combinant des
    arbres de décision entraînés séquentiellement.
  \item[Isolation Forest] Algorithme non supervisé de détection d'anomalies
    par isolation aléatoire des observations.
  \item[Z-score] Écart d'une observation à sa moyenne, exprimé en nombre
    d'écarts-types.
  \item[RFM] \textit{Recency, Frequency, Monetary} : analyse du comportement
    d'achat selon la récence, la fréquence et le montant.
  \item[CLV] \textit{Customer Lifetime Value} : valeur client estimée sur la
    durée de la relation commerciale.
  \item[Churn] Attrition : perte d'un client devenu inactif.
  \item[DAX] \textit{Data Analysis Expressions} : langage de formules de
    Power BI.
  \item[API REST] Interface de programmation exposant des ressources via HTTP,
    ici au format JSON.
  \item[DSF] \textit{Dubai Shopping Festival} : période commerciale majeure
    aux Émirats (décembre--janvier).
  \item[White Friday] Équivalent du \textit{Black Friday} au Moyen-Orient.
\end{description}

% ============================================================================
% CHAPITRE 1 - INTRODUCTION
% ============================================================================
\chapter{Introduction}

\section{Présentation personnelle et contexte du stage}

Je m'appelle BEN SAID Nassim, élève-ingénieur en cinquième et dernière année du
cycle ICY à l'INSA Hauts-de-France. Dans le cadre de mon projet de fin
d'études, j'ai effectué un stage de six mois au sein du cabinet de conseil
Bahgat IT Expert, dont le siège est à Dubaï (Émirats Arabes Unis). Le stage
s'est déroulé intégralement à distance depuis le Maroc, en lien quotidien avec
l'équipe du cabinet.

Après un précédent stage orienté cybersécurité et supervision réseau chez RATP
Connect, je souhaitais consacrer mon stage de fin d'études à la science des
données appliquée aux enjeux métier : la valorisation de la donnée est
aujourd'hui l'un des leviers de compétitivité les plus puissants pour les
entreprises, et le conseil offre un terrain idéal pour en mesurer l'impact
concret. Bahgat IT Expert, cabinet pluridisciplinaire positionné à la croisée
de l'expertise technique et juridique, développe une activité croissante
d'analyse de données et de tableaux de bord décisionnels pour ses clients : ce
contexte m'a offert l'opportunité de concevoir une plateforme complète, de la
donnée brute jusqu'à la restitution pour les décideurs.

\section{Objectifs du rapport}

Ce rapport vise à :
\begin{itemize}
  \item présenter l'entreprise d'accueil, son organisation et mon
    positionnement au sein de ses équipes ;
  \item décrire la problématique métier confiée et la démarche d'ingénierie
    adoptée pour y répondre ;
  \item détailler la conception et la réalisation technique de la plateforme
    (pipeline de données, machine learning, restitution) ;
  \item analyser les résultats obtenus, les difficultés rencontrées et les
    compétences développées ;
  \item dresser un bilan global de cette expérience et de son impact sur mon
    projet professionnel.
\end{itemize}

\section{Plan du rapport}

Le chapitre 2 présente l'entreprise Bahgat IT Expert. Le chapitre 3 expose le
contexte, les objectifs et la méthodologie du projet. Le chapitre 4 décrit la
conception de la plateforme (architecture, choix technologiques, modèle de
données). Les chapitres 5, 6 et 7 détaillent la réalisation : le pipeline de
données, les modules de machine learning, puis les couches de restitution
(application web et Power BI). Le chapitre 8 dresse le bilan du projet et des
compétences acquises, avant la conclusion au chapitre 9.

% ============================================================================
% CHAPITRE 2 - PRÉSENTATION DE L'ENTREPRISE
% ============================================================================
\chapter{Présentation de l'entreprise : Bahgat IT Expert}

\section{L'entreprise}

Bahgat IT Expert est un cabinet de conseil et d'arbitrage spécialisé dans les
technologies de l'information, dont le siège est situé à Dubaï (Émirats Arabes
Unis), avec des bureaux en Arabie Saoudite, à Bahreïn et au Canada. Fort de
plus de trente ans d'expérience, le cabinet est certifié auprès du Ministère de
la Justice des Émirats et des cours fédérales, ce qui lui confère un
positionnement original à la croisée de l'expertise technique et de l'expertise
judiciaire.

Ses activités couvrent un large spectre :
\begin{itemize}
  \item l'investigation numérique (\textit{digital forensics}) et l'analyse
    post-incident ;
  \item la cybersécurité et l'analyse de logiciels malveillants ;
  \item la récupération de données et la préservation de preuves numériques ;
  \item la conception et l'optimisation d'infrastructures IT ;
  \item \textbf{l'analyse de données et la Business Intelligence}, notamment la
    construction de tableaux de bord Power BI pour les décideurs ;
  \item l'accompagnement à l'adoption de l'intelligence artificielle et le
    conseil blockchain ;
  \item l'arbitrage et l'expertise judiciaire dans les litiges à composante
    technologique.
\end{itemize}

Le cabinet intervient pour des clients de secteurs variés : distribution et
e-commerce, construction et bâtiments intelligents, pétrole et gaz,
hôtellerie, finance et fintech. C'est au sein de la practice Data \& Business
Intelligence que s'est inscrit mon stage.

% --- Figure : chiffres clés / présentation de l'entreprise ------------------
\begin{figure}[htbp]
  \includegraphics[width=0.9\textwidth]{assets_rs/chiffres_cles_bahgat.png}
  \caption{Chiffres clés de Bahgat IT Expert}
  \label{fig:chiffres-bahgat}
\end{figure}

La position de Dubaï comme hub technologique régional confère au cabinet un
environnement particulièrement dynamique : les stratégies nationales de
transformation numérique (\textit{UAE Digital Government Strategy}, adoption
massive de l'IA dans les services) tirent une forte demande de conseil en
valorisation de la donnée. Dans ce contexte, la practice Data \& BI du cabinet
cherche à industrialiser son offre : disposer d'une plateforme de référence,
réutilisable d'un client à l'autre, plutôt que de reconstruire des tableaux de
bord ad hoc à chaque mission. Mon projet de fin d'études s'inscrit directement
dans cette ambition.

\section{Les collaborateurs}

Le cabinet est organisé autour de plusieurs pôles d'expertise : investigation
numérique et forensics, cybersécurité, infrastructure, arbitrage et expertise
judiciaire, et enfin data \& intelligence artificielle. Pendant mon stage, j'ai
intégré l'équipe Data \& BI, qui accompagne les clients dans la structuration
de leurs données, la mise en place d'indicateurs et la construction de tableaux
de bord décisionnels.

J'ai travaillé sous la responsabilité directe de mon tuteur de stage,
M.~BAHGAT Ahmed, et en interaction
régulière avec les consultants des autres pôles -- notamment cybersécurité pour
les questions de protection des données, un aspect indissociable des projets
data dans la région.

\section{L'environnement de travail}

Le cabinet est installé dans le quartier de Deira, à Dubaï ; mon stage s'est
toutefois déroulé \textbf{entièrement à distance depuis le Maroc}, une
organisation facilitée par un décalage horaire réduit (environ trois heures)
avec les Émirats. Le rythme de travail reposait sur des rituels clairs : un
point de suivi hebdomadaire avec mon tuteur, une démonstration à chaque jalon
du projet, et des échanges quotidiens avec l'équipe via Microsoft Teams. Cette
organisation a exigé -- et fortement développé -- l'autonomie, la
communication écrite structurée et la documentation systématique des avancées,
la taille humaine des équipes favorisant par ailleurs des échanges directs
avec les consultants seniors.

Sur le plan technique, mon environnement de travail s'articulait autour des
outils suivants :
\begin{itemize}
  \item \textbf{Python} (pandas, NumPy, scikit-learn, statsmodels, Flask) pour
    l'ensemble de la chaîne de traitement et l'application web ;
  \item \textbf{Visual Studio Code} et \textbf{Git} pour le développement et le
    versionnement ;
  \item \textbf{Power BI Desktop} pour la restitution décisionnelle ;
  \item \textbf{SQLite} comme entrepôt de données local de démonstration ;
  \item \textbf{Microsoft Teams} et \textbf{Outlook} pour la communication et
    le suivi d'avancement avec mon tuteur et les équipes.
\end{itemize}

% --- Figure : environnement de travail --------------------------------------
\begin{figure}[htbp]
 \centering
  \includegraphics[width=0.6\textwidth]{assets_rs/oo.jpg}
  \caption{Environnement de travail chez Bahgat IT Expert}
  \label{fig:locaux}
\end{figure}

% ============================================================================
% CHAPITRE 3 - CONTEXTE ET OBJECTIFS DU PROJET
% ============================================================================
\chapter{Contexte, objectifs et méthodologie du projet}

\section{Problématique métier}

Les clients de la practice Data \& BI partagent un constat récurrent : leurs
données sont abondantes mais peu exploitées. Elles sont dispersées dans des
silos hétérogènes -- systèmes d'encaissement (POS), CRM, ERP, plateformes
marketing, analytics web -- et le reporting repose sur des extractions
manuelles consolidées dans des tableurs, avec plusieurs conséquences :

\begin{itemize}
  \item un \textbf{reporting lent et rétrospectif} : les indicateurs sont
    produits mensuellement, souvent avec plusieurs semaines de décalage ;
  \item une \textbf{fiabilité incertaine} : les consolidations manuelles
    introduisent doublons, erreurs de format et incohérences ;
  \item une \textbf{absence de vision prospective} : aucune prévision
    structurée n'alimente les décisions d'approvisionnement ou de
    planification ;
  \item une \textbf{détection tardive des incidents} : erreurs de prix, pannes
    en magasin, fraudes aux retours ne sont découvertes qu'au moment du
    reporting, parfois des semaines après les faits ;
  \item une \textbf{connaissance client limitée} : pas de segmentation
    comportementale permettant de cibler les actions marketing.
\end{itemize}

La mission qui m'a été confiée répond à ce constat : concevoir et développer
une plateforme de Business Intelligence \og pilotée par l'IA \fg{},
centralisée, automatisée et démontrable, que le cabinet puisse présenter à ses
clients et adapter à leurs contextes.

\section{Objectifs et périmètre}

Le cahier des charges établi avec mon tuteur fixait six objectifs :

\begin{enumerate}
  \item collecter et traiter des données issues de \textbf{sources multiples}
    (ventes, retours, clients, produits, marketing, web, stocks) ;
  \item mettre en place des \textbf{pipelines automatisés} d'intégration et de
    fiabilisation des données, avec traçabilité complète des corrections ;
  \item construire un \textbf{entrepôt de données} structuré, performant et
    requêtable par les outils de restitution ;
  \item développer des \textbf{modèles de machine learning} pour la prévision
    du chiffre d'affaires, la détection d'anomalies et la segmentation
    client ;
  \item restituer les analyses dans des \textbf{tableaux de bord interactifs} :
    une application web moderne (Flask) et un rapport \textbf{Power BI}
    multi-pages ;
  \item documenter l'ensemble pour permettre la reprise du projet par les
    consultants du cabinet.
\end{enumerate}

\section{La contrainte de confidentialité et la stratégie de données
  synthétiques}

Les données des clients du cabinet sont couvertes par des accords de
confidentialité stricts : il était exclu de les utiliser pour un projet de
démonstration, et a fortiori de les présenter dans ce rapport ou en soutenance.

Plutôt que de subir cette contrainte, j'en ai fait un atout méthodologique en
développant un \textbf{générateur de données synthétiques haute-fidélité}
simulant trois années d'activité d'un distributeur omnicanal du Golfe --
\og Al Noor Retail Group \fg{}, client fictif -- avec 18 points de vente aux
Émirats, en Arabie Saoudite et à Bahreïn, 320 produits et 40\,000 clients
fidélité, pour environ un million de lignes de vente.

Cette approche apporte un avantage décisif : les données étant générées, la
\textbf{vérité terrain est connue}. Des incidents opérationnels (pannes,
erreurs de prix, fraude aux retours) et des défauts de qualité de données
(doublons, formats hétérogènes) ont été volontairement injectés à des dates et
des ampleurs connues. Chaque brique de la plateforme peut donc être évaluée
\textbf{objectivement} : le pipeline a-t-il corrigé tous les défauts ? Le
détecteur d'anomalies a-t-il retrouvé tous les incidents ? Cette démarche
d'évaluation, impossible avec des données réelles non annotées, constitue l'un
des apports méthodologiques majeurs du projet. La plateforme reste par
ailleurs totalement agnostique de la source : brancher les extractions réelles
d'un client ne modifie aucun composant en aval.

\section{Méthodologie et organisation du projet}

Le projet a été conduit sur six mois selon une démarche itérative inspirée des
méthodes agiles, avec des points d'avancement hebdomadaires avec mon tuteur et
des démonstrations à chaque jalon :

\begin{table}[htbp]
  \centering
  \small
  \begin{tabular}{@{}llp{7.6cm}@{}}
    \toprule
    \textbf{Phase} & \textbf{Période} & \textbf{Contenu} \\
    \midrule
    1. Cadrage        & Mois 1 & État de l'art BI/ML, benchmark des outils,
                                 cahier des charges, maquettes \\
    2. Données        & Mois 2 & Générateur de données synthétiques, modèle en
                                 étoile, pipeline ETL et qualité de données \\
    3. Machine learning & Mois 3--4 & Prévision, détection d'anomalies,
                                 segmentation ; backtests et validation \\
    4. Restitution    & Mois 4--5 & Application web Flask, rapport Power BI,
                                 API, design system \\
    5. Consolidation  & Mois 6 & Tests de bout en bout, documentation,
                                 transfert de compétences, préparation de la
                                 démonstration finale \\
    \bottomrule
  \end{tabular}
  \caption{Découpage du projet en phases}
  \label{tab:phases}
\end{table}

% --- Figure : planning / Gantt ----------------------------------------------
\begin{figure}[htbp]
\centering
   \includegraphics[width=\textwidth]{assets_rs/gant.png}
  \caption{Planning prévisionnel et réalisé du projet}
  \label{fig:gantt}
\end{figure}

% ============================================================================
% CHAPITRE 4 - CONCEPTION
% ============================================================================
\chapter{Conception de la plateforme}

\section{Architecture générale}

La plateforme, baptisée \textbf{Bahgat Insight Platform}, suit une architecture
en couches classique des systèmes décisionnels, orchestrée de bout en bout par
un script unique (\texttt{run\_pipeline.py}) :

\begin{enumerate}
  \item \textbf{Sources} : neuf extractions CSV simulant les systèmes sources
    (POS, CRM, ERP, marketing, web analytics) ;
  \item \textbf{Ingestion et qualité} : lecture, dédoublonnage, normalisation,
    réparation et mise en quarantaine, avec journal d'audit ;
  \item \textbf{Entrepôt} : schéma en étoile (dimensions et faits) et tables
    d'agrégats pré-calculées, stockés dans SQLite ;
  \item \textbf{Machine learning} : trois modules lisant l'entrepôt et y
    réinscrivant leurs résultats (prévisions, alertes, segments) ;
  \item \textbf{Restitution} : application web Flask interrogeant l'entrepôt
    via une API JSON, et export CSV propre vers Power BI.
\end{enumerate}

% --- Figure : architecture --------------------------------------------------
\begin{figure}[htbp]
\centering
   \includegraphics[width=\textwidth]{assets_rs/Schema-archi.png}
  \caption{Architecture générale de la Bahgat Insight Platform}
  \label{fig:architecture}
\end{figure}

\section{Choix technologiques}

Chaque brique a fait l'objet d'un arbitrage entre plusieurs alternatives,
résumé dans le tableau~\ref{tab:choix-techno}.

\begin{table}[htbp]
  \centering
  \small
  \begin{tabular}{@{}lp{4.1cm}p{6.2cm}@{}}
    \toprule
    \textbf{Besoin} & \textbf{Choix retenu} & \textbf{Justification} \\
    \midrule
    Traitement de données & Python, pandas, NumPy &
      Standard de fait, écosystème riche, montée en charge suffisante pour
      $\sim$1,3M de lignes \\
    Entrepôt & SQLite &
      Zéro installation, fichier unique, adapté à une démonstration locale ;
      migration PostgreSQL/Synapse = simple changement de connexion \\
    Prévision & statsmodels (Holt-Winters) &
      Référence des méthodes statistiques de séries temporelles,
      interprétable \\
    ML & scikit-learn &
      Gradient Boosting, Isolation Forest, KMeans : robustes, documentés,
      reproductibles \\
    Application web & Flask + ApexCharts &
      Légèreté, API JSON native, exécution locale en une commande ;
      interface moderne inspirée du design system shadcn/ui \\
    Restitution BI & Power BI Desktop &
      Outil de référence de la practice et attendu par les clients du
      cabinet \\
    \bottomrule
  \end{tabular}
  \caption{Principaux choix technologiques et leurs justifications}
  \label{tab:choix-techno}
\end{table}

Deux choix méritent un commentaire. D'abord, la prévision repose sur un duel
\og champion/challenger \fg{} entre une méthode statistique et un modèle de
machine learning, tranché par une évaluation chiffrée plutôt que par un a
priori technologique (section~\ref{sec:forecast}). Ensuite, l'application web
n'utilise volontairement aucun framework front-end lourd : un design system
maison (inspiré de shadcn/ui), des graphiques ApexCharts et du JavaScript
natif suffisent, ce qui simplifie le déploiement chez un client.

\section{Modèle de données : le schéma en étoile}

L'entrepôt suit la modélisation dimensionnelle de Kimball : les mesures
(montants, quantités) sont portées par des tables de faits, reliées à des
dimensions descriptives partagées.

\begin{itemize}
  \item \textbf{Dimensions} : \texttt{dim\_date} (calendrier enrichi :
    week-ends, Ramadan, Aïd, événements commerciaux), \texttt{dim\_store} (18
    magasins, 3 pays), \texttt{dim\_product} (320 produits, 8 catégories),
    \texttt{dim\_customer} (40\,000 clients), \texttt{dim\_supplier} ;
  \item \textbf{Faits} : \texttt{fact\_sales} ($\sim$1M de lignes au grain
    ligne de ticket), \texttt{fact\_returns}, \texttt{fact\_marketing},
    \texttt{fact\_web\_traffic}, \texttt{fact\_inventory} ;
  \item \textbf{Agrégats} : quatre tables pré-calculées (KPIs quotidiens,
    ventilations par magasin et par catégorie) garantissant des temps de
    réponse inférieurs à 100~ms sur les tableaux de bord ;
  \item \textbf{Sorties ML} : prévisions, alertes, segments et tables de
    validation, réinjectées dans l'entrepôt pour être consommées par les deux
    couches de restitution.
\end{itemize}

% --- Figure : schéma en étoile ----------------------------------------------
\begin{figure}[htbp]
\includegraphics[width=\textwidth]{assets_rs/schema_etoile.png}
  \caption{Schéma en étoile de l'entrepôt de données}
  \label{fig:etoile}
\end{figure}

% ============================================================================
% CHAPITRE 5 - PIPELINE DE DONNÉES
% ============================================================================
\chapter{Réalisation : le pipeline de données}

\section{Génération du jeu de données synthétique}

Le générateur (module \texttt{src/data\_generation}) produit un jeu de données
réaliste et \textbf{reproductible} (graine aléatoire fixée). Le réalisme repose
sur la superposition de plusieurs couches de signal, calibrées sur les
spécificités du retail du Golfe :

\begin{itemize}
  \item une \textbf{saisonnalité hebdomadaire} (pics de fréquentation du
    week-end) et \textbf{annuelle} : Ramadan et Aïd (avec déplacement des
    dates d'une année sur l'autre), Dubai Shopping Festival, White Friday,
    rentrée scolaire, fête nationale ;
  \item une \textbf{tendance de croissance} différenciée : $+9\,\%$ par an pour
    les magasins physiques, $+27\,\%$ pour les deux sites e-commerce ;
  \item des \textbf{promotions} corrélées aux événements commerciaux, avec des
    taux de remise réalistes ;
  \item une \textbf{popularité des produits décroissante avec le prix}
    (loi de type Zipf), propriété essentielle pour obtenir des chiffres
    d'affaires journaliers stables par magasin (cf. section~\ref{sec:anomalies}) ;
  \item un comportement client hétérogène (propension d'achat log-normale),
    condition nécessaire à une segmentation RFM signifiante.
\end{itemize}

Deux familles d'éléments sont ensuite injectées volontairement, constituant la
\textbf{vérité terrain} de l'évaluation (tableau~\ref{tab:incidents}) :

\begin{table}[htbp]
  \centering
  \small
  \begin{tabular}{@{}p{5.6cm}llp{2.6cm}@{}}
    \toprule
    \textbf{Incident injecté} & \textbf{Date(s)} & \textbf{Périmètre} &
    \textbf{Effet} \\
    \midrule
    Panne du système d'encaissement & 14/05/2024 & Deira City Centre &
      CA $\times\,0{,}10$ \\
    Coupure électrique & 17/09/2024 & Riyadh Park & CA $\times\,0{,}25$ \\
    Inondation d'entrepôt & 10--12/02/2025 & 2 magasins Sharjah &
      CA $\times\,0{,}35$ \\
    Méga vente flash & 15/06/2025 & Tous magasins & CA $\times\,1{,}75$ \\
    Bug de prix (électronique $-50\,\%$) & 03/09/2025 & Tous magasins &
      Marge effondrée \\
    Pic e-commerce White Friday & 28--29/11/2025 & 2 sites en ligne &
      CA $\times\,2{,}10$ \\
    Fraude aux retours & 08--14/01/2026 & Sahara Centre &
      Remboursements $\times\,10$ \\
    Rupture de stock électronique & oct. 2025 & Magasin amiral &
      Stocks à zéro \\
    \bottomrule
  \end{tabular}
  \caption{Incidents opérationnels injectés dans les données (vérité terrain)}
  \label{tab:incidents}
\end{table}

S'y ajoutent des \textbf{défauts de qualité de données} représentatifs des
extractions réelles : doublons d'export, formats de dates mélangés (ISO et
JJ/MM/AAAA), identifiants clients manquants, quantités négatives, prix nuls,
libellés de moyens de paiement incohérents (casse, espaces), doublons
d'adresses e-mail.

Au total, le générateur produit neuf fichiers CSV pour environ
1\,032\,000 lignes de ventes brutes % MAJ : sortie console de votre exécution
en une quinzaine de secondes.

\section{Le pipeline ETL et la qualité des données}

Le pipeline ETL (module \texttt{src/etl}) enchaîne trois étapes instrumentées
-- chaque exécution journalise sa durée, ses volumes et chacune de ses
corrections dans des tables dédiées de l'entrepôt.

\textbf{Extraction} : les neuf fichiers sont lus en typage
\og chaîne \fg{} strict, afin qu'aucune conversion implicite ne masque un
défaut de qualité.

\textbf{Transformation} : chaque défaut identifié déclenche une action
adaptée -- suppression (doublons), normalisation (dates, libellés),
réparation par référentiel (prix nuls corrigés depuis le référentiel
produits), ou mise en quarantaine (quantités négatives, retours orphelins).
Le tableau~\ref{tab:dq} présente le bilan réel d'une exécution.

\begin{table}[htbp]
  \centering
  \small
  % MAJ : reprendre les chiffres exacts de VOTRE dernière exécution
  % (page « Data Pipeline » de l'application ou sortie console de l'ETL)
  \begin{tabular}{@{}lrl@{}}
    \toprule
    \textbf{Défaut détecté} & \textbf{Lignes} & \textbf{Action} \\
    \midrule
    Doublons de transactions        & 2\,573  & Supprimés (premier conservé) \\
    Dates au format JJ/MM/AAAA      & 61\,712 & Normalisées en ISO \\
    Quantités négatives ou nulles   & 180     & Mises en quarantaine \\
    Prix nuls ou invalides          & 150     & Réparés via le référentiel \\
    Libellés de paiement incohérents & 41\,128 & Normalisés \\
    Identifiants client manquants   & 252\,978 & Convertis en NULL (invité) \\
    Casse incohérente des villes    & 1\,600  & Normalisée \\
    E-mails clients dupliqués       & 60      & Marqués pour revue \\
    Retours orphelins               & 5       & Mis en quarantaine \\
    \bottomrule
  \end{tabular}
  \caption{Bilan de fiabilisation d'une exécution du pipeline ETL}
  \label{tab:dq}
\end{table}

Les mesures métier sont ensuite dérivées ligne à ligne (chiffre d'affaires
brut, remises, net, coût, marge), puis les dimensions et faits du schéma en
étoile sont construits.

\textbf{Chargement} : l'entrepôt SQLite est reconstruit avec ses index et ses
quatre tables d'agrégats. L'exécution complète du pipeline -- génération,
ETL, machine learning et export Power BI -- s'effectue en un peu plus d'une
minute, % MAJ : durée de votre exécution
ce qui permet de rejouer la chaîne de bout en bout à la demande.

% --- Figure : console pipeline ----------------------------------------------
\begin{figure}[htbp]
\centering
  \includegraphics[width=0.45\textwidth]{assets_rs/etl1.png}
  \includegraphics[width=0.45\textwidth]{assets_rs/etl2.png}
  \includegraphics[width=0.45\textwidth]{assets_rs/etl3.png}
  \includegraphics[width=0.45\textwidth]{assets_rs/etl4.png}
  
  \caption{Exécution du pipeline : journal de fiabilisation des données}
  \label{fig:console-etl}
\end{figure}

% ============================================================================
% CHAPITRE 6 - MACHINE LEARNING
% ============================================================================
\chapter{Réalisation : les modules de machine learning}

\section{Prévision du chiffre d'affaires}
\label{sec:forecast}

\subsection{Démarche champion/challenger}

Pour prévoir le chiffre d'affaires quotidien à un horizon de 90 jours, deux
modèles concurrents sont entraînés puis départagés par une évaluation
rigoureuse :

\begin{itemize}
  \item \textbf{Holt-Winters} (lissage exponentiel triple, champion
    statistique) : modélise le niveau $\ell_t$, la tendance $b_t$ et la
    saisonnalité hebdomadaire $s_t$ de la série :
    \begin{align}
      \ell_t &= \alpha\,(y_t - s_{t-m}) + (1-\alpha)(\ell_{t-1} + b_{t-1})\\
      b_t &= \beta\,(\ell_t - \ell_{t-1}) + (1-\beta)\,b_{t-1}\\
      s_t &= \gamma\,(y_t - \ell_t) + (1-\gamma)\,s_{t-m}
    \end{align}
    avec $m=7$ (saisonnalité hebdomadaire) et une tendance amortie ;
  \item \textbf{Gradient Boosting} (challenger ML) : régression par arbres sur
    des variables calendaires (jour de semaine, mois, Ramadan, événements
    commerciaux) et des retards de la série (valeurs à $J-1$, $J-7$, $J-14$,
    $J-28$, moyennes mobiles), avec prévision récursive multi-pas.
\end{itemize}

\subsection{Évaluation par backtest}

Les deux modèles sont évalués sur une fenêtre de test de 60 jours jamais vue à
l'entraînement (\textit{holdout}), avec la MAPE comme métrique principale :
\begin{equation}
  \mathrm{MAPE} = \frac{100}{n}\sum_{t=1}^{n}
  \left|\frac{y_t - \hat{y}_t}{y_t}\right|
\end{equation}

\begin{table}[htbp]
  \centering
  % MAJ : reprendre les métriques de VOTRE dernière exécution
  % (page « Forecasting » de l'application, tableau « Model tournament »)
  \begin{tabular}{@{}lrrl@{}}
    \toprule
    \textbf{Modèle} & \textbf{MAPE} & \textbf{RMSE (AED)} & \textbf{Statut} \\
    \midrule
    Holt-Winters      & 8{,}55\,\% & 33\,149 & \textbf{Champion} \\
    Gradient Boosting & 10{,}37\,\% & 43\,871 & Challenger \\
    \bottomrule
  \end{tabular}
  \caption{Backtest des deux modèles de prévision sur 60 jours de holdout}
  \label{tab:backtest}
\end{table}

Le modèle statistique l'emporte, ce qui s'explique : la série présente une
saisonnalité hebdomadaire forte et stable, configuration dans laquelle les
méthodes de lissage exponentiel restent très compétitives face aux modèles
d'apprentissage, pour un coût de calcul et une interprétabilité bien
meilleurs. Le champion produit alors la prévision officielle à 90 jours,
assortie d'un intervalle de confiance à 95\,\% dont la largeur croît avec
l'horizon. Des prévisions par catégorie de produits (huit séries
supplémentaires) complètent la prévision globale, avec des MAPE comprises
entre 10 et 24\,\% selon la volatilité des catégories.

% --- Figure : page Forecasting ----------------------------------------------
\begin{figure}[htbp]
\centering
  \includegraphics[width=\textwidth]{assets_rs/app_forecast.png}
  \caption{Prévision du chiffre d'affaires à 90 jours dans l'application}
  \label{fig:app-forecast}
\end{figure}

\section{Détection d'anomalies}
\label{sec:anomalies}

\subsection{Deux détecteurs complémentaires}

La détection d'anomalies combine deux approches travaillant à des échelles
différentes :

\begin{itemize}
  \item \textbf{Isolation Forest} (niveau entreprise, multivarié) : chaque
    journée est décrite par un vecteur de KPIs (chiffre d'affaires, commandes,
    panier moyen, taux de remise, taux de marge, remboursements, rapportés à
    leur référence saisonnière). L'algorithme isole les journées atypiques
    \emph{dans leur combinaison} de indicateurs -- il détecte ainsi le bug de
    prix du 03/09/2025, journée dont le chiffre d'affaires était normal mais
    dont la marge s'était effondrée : invisible pour un contrôle univarié ;
  \item \textbf{Z-score glissant} (niveau magasin, univarié) : le chiffre
    d'affaires de chaque magasin est comparé jour par jour à une référence
    glissante ; les retours font l'objet du même contrôle pour détecter les
    abus de remboursement.
\end{itemize}

\subsection{Une difficulté instructive : queues lourdes et robustesse}

La première version du détecteur magasin comparait chaque journée à la moyenne
et à l'écart-type des 28 jours précédents. Or ce détecteur manquait des
effondrements pourtant flagrants ($-75\,\%$ de chiffre d'affaires), tout en
levant plus d'un millier d'alertes non pertinentes. Le diagnostic, mené en
interrogeant directement l'entrepôt, a révélé deux causes :

\begin{enumerate}
  \item la \textbf{saisonnalité hebdomadaire} gonflait l'écart-type de la
    fenêtre (les week-ends côtoyaient les jours creux), noyant le signal ;
    la solution a consisté à comparer chaque journée à \emph{ses homologues
    du même jour de semaine} (huit dernières occurrences), ce qui élimine
    cette variance parasite ;
  \item la \textbf{distribution à queue lourde} des ventes journalières :
    dans la première version du générateur, la popularité des produits était
    indépendante du prix, si bien qu'un article à 10\,000 AED pouvait être le
    plus vendu de sa catégorie -- une seule vente faisait alors varier de
    20\,\% le chiffre d'affaires du jour, rendant l'écart-type instable sur
    de petits échantillons. La correction a été double : rendre la
    popularité décroissante avec le prix dans le générateur (propriété
    conforme au retail réel), et ajouter au détecteur une règle de
    sauvegarde fondée sur l'ampleur relative de l'écart (une journée à moins
    de la moitié de sa référence est une anomalie, quelle que soit la
    variance locale), assortie d'un plancher de variance et de garde-fous en
    volume sur les retours.
\end{enumerate}

Cette itération illustre un enseignement central du projet : un détecteur
statistiquement correct peut échouer sur des données réalistes, et seule une
évaluation contre une vérité terrain permet de s'en apercevoir et de corriger
méthodiquement.

\subsection{Validation contre la vérité terrain}

Après correction, les détecteurs combinés retrouvent \textbf{la totalité des
incidents injectés} (rappel de 100\,\% sur les sept événements de ventes et
retours), % MAJ : vérifier « 7/7 » sur la page Anomalies après votre
          % dernière exécution complète
avec un volume d'alertes maîtrisé, hiérarchisé en quatre niveaux de sévérité.
La granularité de la vérité terrain permet en outre d'attribuer chaque
détection à la bonne méthode : les incidents localisés (panne, inondation)
sont capturés par le z-score magasin, les incidents systémiques (bug de prix,
vente flash) par l'Isolation Forest.

% --- Figure : page Anomalies ------------------------------------------------
\begin{figure}[htbp]
\centering
 \includegraphics[width=\textwidth]{assets_rs/app_anomaly.png}
  \caption{Détection d'anomalies et validation contre la vérité terrain}
  \label{fig:app-anomalies}
\end{figure}

\section{Segmentation client et risque d'attrition}

\subsection{Méthode}

La segmentation repose sur l'analyse RFM des 37\,000 clients actifs :
% MAJ : nombre exact affiché par le module de segmentation
récence (jours depuis le dernier achat), fréquence (nombre de commandes) et
montant (chiffre d'affaires cumulé). Après transformation logarithmique
(distributions fortement asymétriques) et standardisation, un partitionnement
KMeans en $k=5$ groupes est appliqué. Chaque groupe est ensuite étiqueté
automatiquement à partir de son centroïde (rang de valeur et nuance de
récence), produisant des segments directement parlants pour le métier.

Deux scores individuels complètent la segmentation :
\begin{itemize}
  \item un \textbf{risque d'attrition} par courbe logistique centrée sur le
    seuil d'inactivité de 120 jours :
    $\;p_{\text{churn}} = \left(1 + e^{-(r - 120)/30}\right)^{-1}$
    où $r$ est la récence en jours ;
  \item une \textbf{valeur client annualisée} (CLV) estimée à partir du rythme
    de dépense observé, écrêtée au percentile 99.
\end{itemize}

\subsection{Résultats}

\begin{table}[htbp]
  \centering
  % MAJ : reprendre les chiffres de VOTRE dernière exécution
  % (sortie console du module de segmentation ou page « Customers »)
  \begin{tabular}{@{}lrr@{}}
    \toprule
    \textbf{Segment} & \textbf{Clients} & \textbf{Part du CA} \\
    \midrule
    Champions          & 5\,130  & 45{,}5\,\% \\
    Loyal Customers    & 9\,581  & 34{,}1\,\% \\
    Promising Mid-Tier & 4\,656  & 6{,}3\,\% \\
    Slipping Away      & 11\,271 & 12{,}9\,\% \\
    Hibernating        & 6\,461  & 1{,}2\,\% \\
    \bottomrule
  \end{tabular}
  \caption{Segments clients obtenus par KMeans sur les variables RFM}
  \label{tab:segments}
\end{table}

La structure obtenue est conforme aux régularités du retail : environ 14\,\%
des clients (Champions) concentrent 45\,\% du chiffre d'affaires.
Le croisement segment $\times$ risque d'attrition produit le livrable le plus
directement actionnable de la plateforme : une \textbf{liste de surveillance}
des clients à forte valeur et à risque d'attrition élevé, triée par valeur
annualisée, directement exploitable pour une campagne de rétention ciblée.

% --- Figure : page Customers ------------------------------------------------
\begin{figure}[htbp]
\centering
  \includegraphics[width=\textwidth]{assets_rs/custint.png}
  \caption{Segmentation client et liste de surveillance de l'attrition}
  \label{fig:app-customers}
\end{figure}

% ============================================================================
% CHAPITRE 7 - RESTITUTION
% ============================================================================
\chapter{Réalisation : les couches de restitution}

\section{L'application web Flask}

\subsection{Architecture applicative}

L'application (fichier \texttt{app.py}) suit une architecture simple et
robuste : le serveur Flask expose d'une part six pages HTML, d'autre part une
API JSON d'une vingtaine de points d'accès qui interrogent l'entrepôt SQLite
(avec fenêtres temporelles paramétrables : 30, 90, 365 jours). Le navigateur
construit les graphiques côté client avec la bibliothèque ApexCharts.

Les six pages couvrent l'ensemble des besoins d'analyse :
\begin{enumerate}
  \item \textbf{Executive Overview} : six KPIs avec variation vs période
    précédente, tendance du CA, mix catégories, classement des magasins,
    dernières alertes ML ;
  \item \textbf{Sales Analytics} : performance mensuelle, taux de marge, mix
    canal (magasin vs e-commerce), moyens de paiement, top produits filtrable,
    entonnoir e-commerce ;
  \item \textbf{Forecasting} : prévision à 90 jours avec intervalle de
    confiance, tournoi de modèles, précision par catégorie ;
  \item \textbf{Anomaly Detection} : chronologie des anomalies, journal
    d'alertes, validation contre la vérité terrain ;
  \item \textbf{Customer Intelligence} : segments, démographie, liste de
    surveillance churn ;
  \item \textbf{Data Pipeline} : étapes du pipeline, rapport de qualité de
    données, contenu de l'entrepôt -- la transparence opérationnelle de la
    plateforme.
\end{enumerate}

\subsection{Design system et accessibilité}

L'interface s'inspire du design system shadcn/ui : composants sobres (cartes,
badges, tableaux), typographie système, mode sombre commutable. Les couleurs
des graphiques suivent une palette catégorielle \textbf{validée pour
l'accessibilité} (contrastes et distinguabilité pour les déficiences de vision
des couleurs, ordre des teintes fixe), et les niveaux de sévérité des alertes
associent systématiquement une icône et un libellé à la couleur, afin que
l'information ne repose jamais sur la couleur seule.

% --- Figure : page Overview (clair) -----------------------------------------
\begin{figure}[htbp]
\centering
  \includegraphics[width=\textwidth]{assets_rs/eobw.png}
  \caption{Page Executive Overview de l'application (mode clair)}
  \label{fig:app-overview}
\end{figure}

% --- Figure : page Overview (sombre) + Sales --------------------------------
\begin{figure}[htbp]
\centering
  \includegraphics[width=\textwidth]{assets_rs/eob.png}
  \caption{Mode sombre de l'application}
  \label{fig:app-sombre}
\end{figure}

% --- Figure : page Data Pipeline --------------------------------------------
\begin{figure}[htbp]
\centering
  \includegraphics[width=\textwidth]{assets_rs/datapipeline.png}
  \caption{Page Data Pipeline : transparence du traitement des données}
  \label{fig:app-pipeline}
\end{figure}

\section{Le rapport Power BI}

Le second canal de restitution s'adresse aux utilisateurs de l'écosystème
Microsoft, standard chez les clients du cabinet. Le pipeline exporte
automatiquement dix-huit fichiers CSV propres (schéma en étoile complet et
sorties ML) vers un dossier dédié ; le rapport Power BI est construit sur ces
exports.

Le travail a porté sur trois volets :
\begin{itemize}
  \item le \textbf{modèle sémantique} : relations plusieurs-à-un des faits
    vers les dimensions, table de dates marquée, masquage des clés
    techniques ;
  \item une \textbf{bibliothèque d'une trentaine de mesures DAX} organisée par
    thème (ventes, retours, \textit{time intelligence} avec comparaisons
    année-sur-année et cumuls, marketing, e-commerce, indicateurs issus du
    ML), réutilisable telle quelle sur des données client réelles ;
  \item \textbf{cinq pages de rapport} reprenant la logique de l'application
    web (Overview, Sales, Forecast, Anomalies, Customers), avec un thème JSON
    personnalisé aux couleurs de la plateforme et des segments (slicers)
    synchronisés.
\end{itemize}

% --- Figure : modèle Power BI -----------------------------------------------
\begin{figure}[htbp]
\centering
  \includegraphics[width=\textwidth]{assets_rs/star.png}
  \caption{Modèle sémantique du rapport Power BI}
  \label{fig:pbi-modele}
\end{figure}

% --- Figure : pages Power BI ------------------------------------------------
\begin{figure}[htbp]
\centering
  \includegraphics[width=\textwidth]{assets_rs/pbi1.png}
  \caption{Page Executive Overview du rapport Power BI}
  \label{fig:pbi-overview}
\end{figure}

\begin{figure}[htbp]
\centering
  \includegraphics[width=\textwidth]{assets_rs/pbi3.png}
  \caption{Page d'analyse avancée du rapport Power BI}
  \label{fig:pbi-forecast}
\end{figure}

\section{Automatisation et trajectoire de production}

L'ensemble de la chaîne s'exécute en deux commandes : \texttt{python
run\_pipeline.py} (pipeline complet, environ une minute) puis \texttt{python
app.py} (application web). Chaque étape peut aussi être rejouée isolément
(\texttt{-{}-only etl}, \texttt{-{}-only anomaly}, etc.), ce qui s'est révélé
précieux pendant les itérations de mise au point.

La trajectoire vers un déploiement de production chez un client a été
documentée pour le cabinet : dépôt des extractions dans un stockage objet,
orchestration planifiée (Airflow ou Azure Data Factory), entrepôt PostgreSQL
ou Azure Synapse (le code SQL étant portable), application conteneurisée
derrière une authentification d'entreprise, et publication du rapport sur le
service Power BI avec actualisation planifiée. La modularité de la plateforme
rend cette migration incrémentale : chaque brique peut être remplacée sans
toucher aux autres.

% ============================================================================
% CHAPITRE 8 - BILAN
% ============================================================================
\chapter{Bilan du projet et compétences développées}

\section{Résultats obtenus}

Au terme des six mois, la plateforme répond à l'ensemble du cahier des
charges. Les résultats se résument en quelques chiffres :
% MAJ : actualiser ces chiffres avec votre dernière exécution complète

\begin{itemize}
  \item environ \textbf{1,3 million de lignes} ingérées à chaque exécution,
    dont $\sim$1,03 million de lignes de ventes fiabilisées ;
  \item plus de \textbf{360\,000 corrections ou décisions de qualité de
    données} appliquées et journalisées par exécution ;
  \item une prévision à 90 jours avec une \textbf{MAPE de 8,55\,\%} en
    backtest, le champion statistique battant le challenger ML ;
  \item un \textbf{rappel de 100\,\%} du détecteur d'anomalies sur les
    incidents injectés, avec un volume d'alertes maîtrisé ;
  \item cinq \textbf{segments clients} exploitables et une liste de
    surveillance churn triée par valeur ;
  \item une exécution de bout en bout en \textbf{à peine plus d'une minute},
    reproductible à l'identique (graine fixée) ;
  \item deux couches de restitution complètes (application web de six pages,
    rapport Power BI de cinq pages) partageant le même modèle sémantique.
\end{itemize}

Pour le cabinet, la plateforme constitue désormais un support de démonstration
commerciale et une base de code réutilisable pour les missions data : la
stratégie de données synthétiques permet de la présenter à un prospect sans
exposer la moindre donnée réelle.

\section{Difficultés rencontrées et solutions apportées}

\begin{itemize}
  \item \textbf{Robustesse statistique face aux queues lourdes} : le cas
    détaillé en section~\ref{sec:anomalies} (détecteur aveuglé par la
    variance de petits échantillons) a été la difficulté technique la plus
    formatrice, résolue par un diagnostic chiffré dans l'entrepôt, une
    correction du réalisme du générateur et des garde-fous de détection ;
  \item \textbf{Rigueur d'évaluation sans fuite de données} : la construction
    des backtests (fenêtre de holdout jamais vue, références glissantes
    strictement antérieures au jour évalué) a demandé une vigilance
    constante pour garantir des métriques honnêtes ;
  \item \textbf{Hétérogénéité des données simulées} : reproduire des défauts
    réalistes (formats de dates mélangés, doublons d'export) et les corriger
    proprement a représenté un vrai travail d'ingénierie de la donnée ;
  \item \textbf{Environnement Windows} : quelques spécificités de
    l'écosystème Python sous Windows 11 (détection des c\oe urs processeur
    par les bibliothèques de parallélisme, avertissements de convergence des
    optimiseurs) ont nécessité des correctifs de configuration pour obtenir
    une exécution silencieuse et démontrable.
\end{itemize}

\section{Compétences acquises}

\textbf{Compétences techniques} :
\begin{itemize}
  \item ingénierie des données : modélisation dimensionnelle (Kimball),
    conception de pipelines ETL instrumentés, stratégie de qualité de
    données avec audit ;
  \item science des données : séries temporelles (lissage exponentiel,
    validation par backtest), détection d'anomalies supervisée par vérité
    terrain, clustering et scoring client ;
  \item développement full-stack : API REST Flask, front-end JavaScript,
    design system et accessibilité des visualisations ;
  \item Business Intelligence : modèle sémantique Power BI, langage DAX,
    conception de rapports multi-pages ;
  \item bonnes pratiques : code modulaire et documenté, reproductibilité,
    versionnement Git.
\end{itemize}

\textbf{Compétences organisationnelles et humaines} :
\begin{itemize}
  \item conduite d'un projet de six mois en autonomie, du cadrage à la
    livraison, avec des jalons de démonstration réguliers ;
  \item communication avec des interlocuteurs non techniques : vulgarisation
    des choix de modèles et des métriques auprès des consultants ;
  \item travail à distance dans un contexte international et multiculturel :
    communication écrite structurée, gestion du décalage horaire, anglais au
    quotidien.
\end{itemize}

\section{Compétences à approfondir}

Ce projet a également révélé des axes de progression que je souhaite
poursuivre : l'industrialisation à grande échelle des pipelines
(orchestrateurs type Airflow, infrastructures cloud), les architectures MLOps
(suivi de modèles en production, réentraînement automatique, détection de
dérive), et l'approfondissement des modèles de prévision multivariés intégrant
des variables exogènes (promotions, météo, marketing).

\section{Limites et perspectives}

La plateforme, conçue comme démonstrateur, présente des limites assumées qui
dessinent sa feuille de route : entrepôt local à migrer vers une base
serveur pour la production, prévision univariée à enrichir de covariables,
score d'attrition heuristique à remplacer par un modèle supervisé lorsque des
étiquettes de churn réelles seront disponibles, boucle de retour analyste
(confirmation/rejet des alertes) à intégrer au détecteur d'anomalies, et
actualisation planifiée du rapport Power BI via le service en ligne. Chacune
de ces évolutions a été documentée à destination des équipes du cabinet.

% ============================================================================
% CHAPITRE 9 - CONCLUSION
% ============================================================================
\chapter{Conclusion et synthèse de l'expérience}

Ce stage de fin d'études chez Bahgat IT Expert m'a permis de concevoir et de
livrer, de bout en bout, une plateforme de Business Intelligence pilotée par
l'intelligence artificielle : pipeline de données automatisé et audité,
entrepôt dimensionnel, prévision évaluée par backtest, détection d'anomalies
validée contre une vérité terrain, segmentation client actionnable, et deux
couches de restitution professionnelles. Au-delà de la réalisation technique,
c'est la démarche d'ingénierie qui constitue à mes yeux l'apport principal de
cette expérience : formuler des hypothèses, instrumenter, mesurer, diagnostiquer
et corriger -- y compris lorsque le défaut se situe dans les données plutôt que
dans le code.

Travailler au sein d'un cabinet de conseil m'a par ailleurs confronté à une
exigence complémentaire de la technique : la capacité à rendre les résultats
intelligibles et convaincants pour des décideurs, qui conditionne l'impact
réel d'un projet data. La contrainte de confidentialité, transformée en
stratégie de données synthétiques avec vérité terrain, illustre cette
articulation entre rigueur d'ingénieur et sens du métier.

Cette expérience internationale -- menée à distance depuis le Maroc pour un
cabinet de Dubaï, dans une région en pleine accélération sur l'intelligence
artificielle -- confirme mon projet
professionnel : évoluer dans l'ingénierie de la donnée et ses applications
décisionnelles, à l'interface entre la technique et les enjeux métier. Les
compétences consolidées durant ces six mois -- et les axes de progression
identifiés -- constituent une base solide pour débuter ma carrière
d'ingénieur.

Je tiens à remercier une nouvelle fois l'équipe de Bahgat IT Expert et mes
tuteurs pour leur accompagnement tout au long de cette expérience
particulièrement enrichissante, qui marque l'aboutissement de ma formation à
l'INSA Hauts-de-France.

% ============================================================================
% BIBLIOGRAPHIE
% ============================================================================
\begin{thebibliography}{10}
\addcontentsline{toc}{chapter}{Bibliographie}

\bibitem{kimball} R. Kimball, M. Ross, \textit{The Data Warehouse Toolkit:
  The Definitive Guide to Dimensional Modeling}, 3\ieme{} édition, Wiley,
  2013.

\bibitem{hyndman} R.~J. Hyndman, G. Athanasopoulos, \textit{Forecasting:
  Principles and Practice}, 3\ieme{} édition, OTexts, 2021.
  \url{https://otexts.com/fpp3/}

\bibitem{isoforest} F.~T. Liu, K.~M. Ting, Z.-H. Zhou, \og Isolation
  Forest \fg{}, \textit{IEEE International Conference on Data Mining}, 2008.

\bibitem{sklearn} F. Pedregosa et al., \og Scikit-learn: Machine Learning in
  Python \fg{}, \textit{Journal of Machine Learning Research}, vol.~12, 2011.
  \url{https://scikit-learn.org}

\bibitem{statsmodels} S. Seabold, J. Perktold, \og Statsmodels: Econometric
  and Statistical Modeling with Python \fg{}, \textit{Proceedings of the 9th
  Python in Science Conference}, 2010. \url{https://www.statsmodels.org}

\bibitem{rfm} A.~J. Christy et al., \og RFM ranking -- An effective approach
  to customer segmentation \fg{}, \textit{Journal of King Saud University --
  Computer and Information Sciences}, 2021.

\bibitem{flask} Documentation officielle Flask.
  \url{https://flask.palletsprojects.com}

\bibitem{powerbi} Microsoft, Documentation Power BI et référence DAX.
  \url{https://learn.microsoft.com/power-bi/}

\bibitem{pandas} W. McKinney, \og Data Structures for Statistical Computing in
  Python \fg{}, \textit{Proceedings of the 9th Python in Science Conference},
  2010. \url{https://pandas.pydata.org}

\bibitem{bahgat} Site institutionnel de Bahgat IT Expert.
  \url{https://www.bahgatexpert.com}

\end{thebibliography}

% ============================================================================
% ANNEXES
% ============================================================================
\appendix

\chapter{Structure du dépôt de code}

\begin{lstlisting}[language={}, caption={Arborescence du projet}]
Internship project/
|-- run_pipeline.py            # orchestrateur (gen -> etl -> ml -> export)
|-- app.py                     # application web Flask
|-- requirements.txt
|-- src/
|   |-- config.py              # chemins et constantes
|   |-- data_generation/generate_data.py
|   |-- etl/pipeline.py
|   |-- ml/forecasting.py
|   |-- ml/anomaly_detection.py
|   |-- ml/segmentation.py
|   |-- powerbi_export.py
|-- templates/  static/        # interface web (6 pages)
|-- powerbi/                   # theme, mesures DAX, guide, exports CSV
|-- data/                      # brut + entrepot SQLite (genere)
|-- docs/                      # architecture, dictionnaire, guides
\end{lstlisting}

\chapter{Extrait du dictionnaire de données}

Cette annexe synthétise le dictionnaire de données de la plateforme :
extractions sources, tables de l'entrepôt (grain, colonnes, volumétrie) et
sorties des modules de machine learning.

\begin{figure}[htbp]
  \centering
  \includegraphics[width=\textwidth]{assets_rs/dico_sources.png}
  \caption{Dictionnaire de données — extractions sources}
\end{figure}

\begin{figure}[htbp]
  \centering
  \includegraphics[width=\textwidth]{assets_rs/dico_entrepot.png}
  \caption{Dictionnaire de données — entrepôt (dimensions, faits, agrégats)}
\end{figure}

\begin{figure}[htbp]
  \centering
  \includegraphics[width=\textwidth]{assets_rs/dico_ml.png}
  \caption{Dictionnaire de données — sorties ML et tables d'exploitation}
\end{figure}

\chapter{Extraits de code significatifs}

\section*{Extrait 1 -- Détection d'anomalies par z-score même-jour-de-semaine}

\begin{lstlisting}[language=Python, caption={Baseline même-jour-de-semaine et
  règle de sauvegarde (extrait de \texttt{src/ml/anomaly\_detection.py})}]
dow = g.index.dayofweek
mu = g.groupby(dow)["v"].transform(
    lambda s: s.rolling(8, min_periods=4).mean().shift(1))
sd = g.groupby(dow)["v"].transform(
    lambda s: s.rolling(8, min_periods=4).std().shift(1))
sd = np.maximum(sd, 0.06 * mu)      # plancher de variance
g["z"] = (g["v"] - mu) / sd
rel = (g["v"] - mu).abs() / mu.replace(0, np.nan)
# regle principale + regle de sauvegarde (effondrement massif)
strong_drop = (g["z"] <= -2.0) & (rel >= 0.50)
hit = g[((g["z"].abs() >= SEUIL_Z) & (rel >= 0.30)) | strong_drop]
\end{lstlisting}

\section*{Extrait 2 -- Mesures DAX de time intelligence}

\begin{lstlisting}[language={}, caption={Mesures DAX (extrait de
  \texttt{powerbi/DAX\_measures.dax})}]
Net Revenue = SUM ( FactSales[net_amount] )

Revenue LY =
CALCULATE ( [Net Revenue], SAMEPERIODLASTYEAR ( DimDate[date] ) )

Revenue YoY % =
DIVIDE ( [Net Revenue] - [Revenue LY], [Revenue LY] )

Return Rate % = DIVIDE ( [Refund Amount], [Net Revenue] )
\end{lstlisting}

\chapter{Captures d'écran complémentaires}

\begin{figure}[htbp]
\centering
  \includegraphics[width=\textwidth]{assets_rs/Screenshot 2026-08-03 092419.png}
  \caption{Page Sales Analytics app.py}
\end{figure}
\begin{figure}[htbp]
\centering
  \includegraphics[width=\textwidth]{assets_rs/pbi2.png}
  \caption{Page Sales Analytics Power BI}
  \end{figure}
\begin{figure}[htbp]
\centering
  \includegraphics[width=\textwidth]{assets_rs/pbi4.png}
  \caption{Page Anomaly Detection Power BI}
  \end{figure}
\begin{figure}[htbp]
\centering
  \includegraphics[width=\textwidth]{assets_rs/pbi5.png}
  \caption{Page Customer Intelligence Power BI}
  \end{figure}
\end{document}
